\subsection{Discussion}
Our results support the paper's main claims across all experiments: registers consistently improve object discovery, absorb outlier token content, and maintain local feature quality in normal patches. The directional effects are robust across model families and datasets.

\subsubsection{What was easy}
The qualitative experiments (``LOST intermediate maps``, ``Norm distributions``, ``Positional focus'') were straightforward to implement and produced results that closely matched the paper visually. The forward hook pattern in PyTorch is well-documented and the timm/open\_clip model APIs are stable, which made feature extraction reliable once the correct hook targets were identified. The paper's descriptions of LOST (``Object discovery``) and the norm distribution experiment were precise enough to implement directly without ambiguity. The linear probing results for global information (Table~\ref{tab:outlier_token_probes}) matched the paper's qualitative conclusions closely with the giant model, and the direction of all effects replicated without exception.

\subsubsection{What was difficult}
\paragraph{LOST Gram matrix.} 
\label{sec:LOST_gram_matrix}
Without per-row mean-centering, DINOv2 keys produce an all-positive similarity matrix causing degree-uniformity collapse — the seed selection degenerates to a fixed corner patch for every image. The paper's formula includes no centering step, and understanding why our setup required it, and at what cost to absolute accuracy, needed multiple debugging iterations. The resulting 15-20\% CorLoc gap could not be fully closed without access to the original training configurations.
\paragraph{Unavailable checkpoints.} Neither DeiT-III+reg nor OpenCLIP+reg weights were released. The OpenCLIP test-time registers proxy \citep{jiang2025vision} is conceptually different from a natively trained model; the DeiT-III proxy uses a different architecture and training paradigm entirely. Both limit the interpretability of those rows in Table~\ref{tab:object_discovery}.
\paragraph{Pre-LayerNorm hook.} Appendix D.1 required hooking the input of the final LayerNorm, not its output — post-LayerNorm activations are rescaled by learned parameters and destroy the bimodal norm signal. This is not stated in the paper and was found through experimentation.
\paragraph{OpenCLIP values hook.} For Figure~\ref{fig:seed_expansion_openclip}, the correct hook target for OpenCLIP values is the full ResidualAttentionBlock output, not the QKV projection. The pre-LN normalisation erases per-patch norm differences before the value projection, making direct V extraction uniformly flat. Identifying this required careful reading of the OpenCLIP architecture.